\documentclass[11pt]{beamer}
\usetheme{metropolis}

\usepackage{fontspec}
\setmainfont{Fira Sans}

%\usepackage{polyglossia}
%\setdefaultlanguage{german}

\usepackage{amsmath, amsfonts, amssymb}

\usepackage{tikz}
\usepackage{graphicx}
\usepackage{caption}
\usepackage{comment}
\usepackage{listings}
\lstset{
%	language=Python,
	basicstyle=\ttfamily,
	breaklines=true,
	numbers=left,
	numberstyle=\tiny,
	frame=tb,
	tabsize=3,
	columns=fixed,
	showstringspaces=false,
	showtabs=false,
	%keepspaces,
	commentstyle=\color{gray},
	keywordstyle=\color{orange}
	}


\author{Leon Rauschning}
\title{Invisible No More}
\subtitle{Common sense visualization with ggplot}
\institute{LMU und TU Munich}
%\date{}

\newcommand{\red}[1]{\textcolor{red}{#1}}
\newcommand{\gray}[1]{\textcolor{gray}{#1}}

\newcommand\blfootnote[1]{%
\begingroup
\renewcommand\thefootnote{}\footnote{#1}%
\addtocounter{footnote}{-1}%
\endgroup
}

\begin{document}
\begin{frame}
	\maketitle
\end{frame}

% brainstorm
\section*{Why do we visualize data?}

% lead into ggplot
% decompose elements of a graphic
% intuition: languages need grammar
% stay on concept level, maybe brainstorm some visualizations in class?
% descriptive, not prescriptive
% language analogy: no snake sentences?


\begin{frame}{}

\end{frame}

\begin{frame}{Grammars of Graphics}

\end{frame}

% next pres
% explain different layers of the grammars
% go into ggplot API 

\end{document}
